\chapter{System Integration and Engineering Implementation\label{cha:chapter7}}

Chapters~5 and~6 describe the theory behind our hybrid battery model, spatial mutual exclusion, AMCL recoveries, and re-allocation strategies. In other words, they answer the question of what our fleet should do and why. However, saying what the system should do is only half of the battle, especially when it comes to something as distributed and multi-layered as ROS~2. How exactly do we turn theoretical claims about allocation policies and claim predicates into actual working Python nodes with launch configurations and software packages that can run on real hardware? This is the question that this chapter attempts to answer~\cite{quigley2015programming}. While most of 5 and 6 are concerned with describing the properties of a certain state space, this chapter is about the nuts and bolts of process graphs, TF trees, and thread boundaries. Put simply, we have to get the engineering fundamentals right before any of the higher-level mathematical machinery can be useful in practice.

The remainder of this chapter is structured as follows. Section~7.1 discusses some of the packaging and topological considerations that went into the final code base, including the rationale behind merging several previously separate classes into one node. Section~7.2 describes the launch file organization, including our transition from a brittle timer-based approach to a more robust gate-based launch strategy enabled by four custom polling nodes that act as readiness sensors. Section~7.3 explains how we mounted three namespaced robots onto a single shared TF tree and the associated challenges of dealing with URDF frame prefixes, transform duplication, and the sole exception in the form of global topics. Section~7.4 discusses the passive diagnostic tools we implemented for the fleet, including a clever dual-subscription technique and timeline reconciliation mechanism that let us tap into the robot state without disrupting ongoing control loops. Section~7.5 describes the calibration tools we developed, namely the odometry covariance relay and yaw-align script. Section~7.6 lifts the lid on the PyQt5 operator interface and makes the case for using WSLg in this setting. Section~7.7 talks about the Gazebo simulation setup and Cartographer-based offline localization and mapping pipeline. Finally, Section~7.8 wraps up the chapter with a few notes on design philosophy, namely the allocation of ownership and responsibilities, explicit notions of instantiation readiness, and failure gradients. With these engineering fundamentals in place, we can finally move on to discussing the results in Chapter~8.

\section{Software Architecture: Package Topology and Process Consolidation}

\subsection{Ament Python Package Layout and the Consolidated Node Graph}

Unlike most larger ROS~2 applications, which tend to split out perception, planning, and coordination logic into separate Ament Python packages for stricter compile-time and dependency isolation~\cite{macenski2022robot}, the autonomy stack is built as a monolithic \texttt{warehouse\_multi\_robot} Ament Python package. The practical benefits of such a layout, in this case, greatly outweigh the costs: with a reduced robot count of three, managing multiple \texttt{package.xml} manifests and importing cross-package symbols becomes a significantly more expensive proposition than it would be in a larger system.

The package itself keeps the resource classes separated by functionality: the \texttt{launch/} directory contains the deterministic launch scaffolding described in Section~7.2, the \texttt{config/} houses the robot-specific Nav2 parameter overlays, bridge YAMLs, and the Cartographer Lua config~\cite{macenski2024modular}; the \texttt{maps/} retains the occupancy grid from the offline SLAM pass of Section~7.7.2, and the \texttt{models/} and \texttt{worlds/} contain the namespaced SDF robot descriptions and the modified warehouse world, respectively. The \texttt{warehouse\_multi\_robot\-/warehouse\_multi\_robot/} subtree contains the Python code itself, following the standard Ament Python convention which isolates the installable Python package contents from the rest of the ROS resource tree.

Rather than listing all the configuration, map, and world files that need to be installed alongside the Python code in \texttt{setup.py} explicitly, the build script traverses the \texttt{launch}, \texttt{config}, \texttt{maps}, \texttt{worlds}, and \texttt{models} directories at build-time, walking every subdirectory it finds and collecting the individual files present within each into the install-set automatically. A manually maintained \texttt{data\_files} list poses significant "works-in-progress" risks: every time a new configuration variant is introduced (e.g., a new Nav2 parameter overlay or a bridge YAML) or an asset is moved around, \texttt{setup.py} needs to be adjusted to reflect that, and an omission results in a file being absent from the installed package. By contrast, a recursive directory walk ensures that whatever files are in the resource directories at the time of the build are properly tracked, each grouped under its own containing directory to match the install-path format \texttt{setuptools} itself expects for a \texttt{data\_files} entry.

On the process level, the package exposes six entry-point defined console-scripts: \texttt{waypoint\_sender}, \texttt{battery\_monitor}, \texttt{mission\_gate}, \texttt{imu\_relay}, \texttt{monitoring\_dashboard}, and \texttt{window\_tiler}; every one of these save for \texttt{window\_tiler} instantiates and spins its own \texttt{rclpy} node. This number is lower than one might expect from a three-robot system with localization, mapping, and navigation functions: in particular, the \texttt{waypoint\_sender} binary serves as the consolidation point for multiple monitoring roles. Internal to this node are the \texttt{FleetHealthMonitor}, the \texttt{ClaimManager}, and the \texttt{TaskCoordinator}, which handle the robot state observation and allocation logic; they are implemented as Python objects alongside the main \texttt{WaypointSenderNode} class. The reasoning behind this arrangement, as well as its practical trade-offs, are discussed in Section~7.1.2.

\subsection{Rationale for Helper-Class Consolidation over Multi-Node Decomposition}

In a more traditional ROS~2 arrangement, the \texttt{FleetHealthMonitor}, the \texttt{ClaimManager}, and the \texttt{Task\-Coordinator} would be implemented as separate nodes. Each of them would have a correspondent Node subclass and would receive updates from the \texttt{waypoint\_sender} node via internal topics. Instead, the three monitors are implemented as regular Python objects, which allows for their internal state to be directly observable by the parent \texttt{WaypointSenderNode} object. This approach was the default for the earlier versions of this package: the \texttt{agent\_coordinator} console script was implemented as a separate node with its own lifecycle, until recently moved into the \texttt{waypoint\_sender} process along with the main coordination logic. This design choice is discussed below from the perspective of the potential gains and drawbacks of this consolidation.

The most immediate advantage of this arrangement stems from the $10\,\text{Hz}$ coordination loop frequency, which imposes strict limits on the allowable overhead for inter-node communication. As the \texttt{select\_next\-\_target} method demonstrates, checks against shared fleet state information (\texttt{active\_claims}, peer battery levels) and updates to the \texttt{completed\_waypoints} set all occur within this control loop: having the shared state represented as regular Python objects saves on the overhead of serialization, deserialization, and transport. In particular, the module docstring for \texttt{waypoint\_sender.py} explicitly states that the decision to consolidate the three helper classes into the single node was motivated by the desire to perform heartbeat checking, task claiming, and re-allocation "locally in memory with zero network latency~\cite{casini2023response}."

This benefit comes at the expense of reduced fault isolation: in this architecture, a fatal exception in one of the helper classes will take down the entire process, whereas in a multi-node arrangement, other ROS~2 nodes would be able to continue operating. In practice, however, such failures would be rare in this particular system: by virtue of being a three-robot system, it was much more straightforward to implement the required level of fault tolerance at the process level for the entire agent rather than for individual nodes. The Chapter~6 considerations on fault tolerance are indeed applicable to the design under review. In particular, the heartbeat timeout of Equation~\ref{eq:heartbeat_timeout_formula} and the battery level check that triggers the transition to the fail-stop state are implemented on the robot level, not the node level. As a result, if any helper class (such as \texttt{ClaimManager}) were to malfunction, it would bring down the entire robot process, and supervision of sub-processes would be of little value, providing no significant benefit over the heartbeat mechanism.

It is also notable that the \texttt{waypoint\_sender} node is not the only process which supervises helper classes: the \texttt{battery\_monitor} and the \texttt{mission\_gate} nodes are separate ROS~2 nodes, each with their own console-script entry point. Both of them implement the critical, but independent, domains of the robot lifecycle: the battery simulation must always be running in order to reliably publish the FAILED transition when the decision logic in the \texttt{waypoint\_sender} node becomes unresponsive. Similarly, the mission-arming logic is explicitly external to the \texttt{waypoint\_sender} node: as the name suggests, it acts as a gate for the robot to engage in coordinated navigation. It is not a node supervisor itself, but it is a logic trigger for the external operator to pause coordinated activity and for individual robots to resume it independently. In general, if the monitored entity or the logic domain has a lifecycle that is loosely coupled to the parent node's coordination loop, it is better to expose it as a separate node~\cite{teper2024timing}. On the other hand, if the helper class lifecycle is fully dependent on the parent node's execution context, it is better to consolidate it into a single process, as the \texttt{waypoint\_sender}/\texttt{agent\_coordinator} exemplifies~\cite{macenski2023impact}.

\section{Deterministic Launch Orchestration}

\subsection{Failure Modes of Timer-Based Sequencing under Variable Boot Latency}

An earlier version of the simulation bring-up procedure, which is preserved in \texttt{slam.launch.py} for standalone mapping runs, arranges gazebo, sensor bridge, robot state publisher, and cartographer nodes in a sequence of \texttt{TimerAction} blocks with hardcoded $5.0\,\text{s}$, $8.0\,\text{s}$, $15.0\,\text{s}$, and $16.0\,\text{s}$ respectively. Each such delay is an estimate of the time needed for the previous element to initialize and publish its topics: five seconds for the physics server to register a spawn request, another three seconds for the sensor bridge and robot-state publisher to bind their subscribers, etc., and so on for the remaining elements and their dependencies. Launch tree sequences constructed from such \texttt{TimerAction}s are common in early ROS~2 tutorials; they are simple to implement and usually effective so long as a developer works within approximately the same hardware environment, with comparable load during testing and in production.

The problem with these hardcoded delays is precisely that each is calibrated to a specific machine with a particular set of background processes, and gazebo's physics engine is especially sensitive to such variations. For example, the time spent in the Gazebo world load phase depends critically on whether Ogre2's shader programs are warm from a previous run, what other GPU-intensive processes are competing for resources at the moment, and the virtualization host's performance, which in turn is affected by the WSL2 environment used in this project. All of these factors affect the reliability of a $5.0\,\text{s}$ delay at the beginning of \texttt{slam.launch.py}, which serves as the allocated window for Gazebo to finish world loading. A too-conservative estimate will waste real time on every launch, while an overly optimistic one may cause downstream elements of the launch tree to fail silently.

Such failures are especially common and difficult to diagnose for processes that have nothing to do with timing, but which happen to be placed by the launcher after a slow operation. The \texttt{ros\_gz\_bridge} node launched by \texttt{slam.launch.py} immediately after Gazebo's world is announced to be loaded can fail to register any topics if Gazebo's topic discovery takes longer than expected to complete. This is a plausible outcome whenever the physics server's own discovery lag exceeds the allotted window: the launch tree proceeds successfully to a later Node action, which runs \texttt{ros\_gz\_bridge}'s \texttt{parameter\_bridge} executable but does not bind to any topics because there are none to bind to, and the only symptom is a diagnostic message to the effect that no matching Gazebo topics were found to bridge.

Generally speaking, a launch tree designed around \texttt{TimerAction} blocks has no inherent mechanism to detect such situations, only a set of assumptions about what operations take how long. As a result, the five hardcoded delays in \texttt{slam.launch.py} encode five separate opportunities for such a situation to occur, and in each case, the resolution is far from obvious. The failure to detect a missing topic is not even a failure in the conventional sense: the \texttt{ros\_gz\_bridge} process is not terminated prematurely, so the launch tree sees all elements as successfully launched, even though nothing is published. Any diagnostic message that would identify this condition as something worth user attention is drowned out by the overall success of the launch tree, which is now erroneously reporting all its nodes as active.

The issue with such non-fatal launch failures is that a \texttt{TimerAction}-based launch tree has no built-in way of handling them. ROS~2 launch library primitives are fundamentally opportunistic about node spawning: \texttt{TimerAction} simply delays execution of its child actions, but says nothing about whether their preconditions are met. As a result, the five hardcoded delays in \texttt{slam.launch.py} each represent an independent potential source of such a precondition violation, which will manifest as a generic "nothing is publishing" error unrelated to its true cause~\cite{blass2021automatic}.

Two additional factors contribute to the particular danger of using such a launch tree structure for multi-robot processes. First, the delays are not additive: attempting to launch a third or a fourth robot using \texttt{slam.launch.py} would either require adjusting every downstream element's waiting time to account for the increased system load or cause unnecessary delays in single-robot runs to accommodate the hypothetical four-robot case. Second, such a launch tree has no way of distinguishing between a process that failed to launch altogether and one that is simply taking longer than expected. The behavior of a Gazebo world that has failed to load due to an incorrect file path is not appreciably different, from the perspective of a \texttt{TimerAction}-based launcher, from a world that needs an extra eight seconds of patience to finish loading. The readiness-polling infrastructure described below was implemented in the production multi-robot launch tree to address these issues.

\subsection{Gate-Based Readiness Polling Primitives}

The production multi-robot launch tree replaces every assumption encoded in a \texttt{TimerAction} block with an explicit readiness gate: a polling subprocess, implemented in \texttt{gate.py}, that waits for an appropriate system signal before allowing a Node action to proceed. The four readiness gates currently implemented in the multi-robot launch tree are waiting for gazebo world load completion, waiting for clock publisher to appear, waiting for topics to be published, and waiting for lifecycle node status. Each is a separate Python script, sharing a common polling infrastructure, which gets invoked by the launch tree to wait for the given condition to be met.

The world readiness gate is the first such check to be triggered, as gazebo world load always precedes the creation of any robot-specific topics. It polls \texttt{gz topic -l} to look for either a \texttt{/world/<name>/stats} or \texttt{/world/<name>/clock} topic matching the world name already derived from the world file's own filename, falling back to a match against the more permissive pattern \texttt{/world/*/(stats\textbar clock)} whenever the exact expected name fails to turn up among the discovered topics. This fallback exists because the caller computes the expected world name ahead of time from the world file path, and a permissive pattern match still confirms that Gazebo has finished loading some world even if that derived name were to disagree with Gazebo's own internal naming for any reason.

The clock topic readiness gate, in contrast, uses \texttt{ros2 topic echo --once /clock --field clock.sec} rather than inspecting the topic list. A \texttt{/clock} topic appearing in the output of \texttt{ros2 topic list} confirms only that a clock publisher has been discovered, but the fact that a publisher is present is not sufficient reason to proceed. Every node that declares \texttt{use\_sim\_time: true} will bind to \texttt{/clock}, but they will all wait for a message to be published before entering active mode, and Nav2's lifecycle nodes will not transition to the next state unless they see periodic updates. The \texttt{--once} argument to \texttt{ros2 topic echo} makes the readiness check wait for a message to appear before proceeding.

The topics readiness gate serves a similar purpose but applies it to an arbitrary set of topics rather than a single \texttt{/clock}. It is implemented using \texttt{ros2 topic info --verbose} and a regular expression match against \texttt{Publisher count} field, blocking until all the requested topics have at least one publisher. The per-robot launch chain, described in Section~7.2.3, will call this script to wait for appearance of the five required topics: laser scan, filtered odometry, robot joints, and the two TF trees, since Nav2's lifecycle nodes are configured to only transition to the active state when all of these are available simultaneously.

The lifecycle readiness gate operates at a higher level and queries \texttt{ros2 lifecycle get <node>} for each node in a specified list, waiting for either \texttt{active} or \texttt{active~[3]} output, depending on the ROS distribution. The presence of the \texttt{active} state in the output is taken as a signal that the node is now fully initialized and ready to accept configuration or mission commands. This gate is currently used to wait for all six nodes managed by Nav2 (\texttt{map\_server}, \texttt{amcl}, \texttt{planner\_server}, \texttt{controller\_server}, \texttt{behavior\_server}, \texttt{bt\_navigator}) before launching the mission-level \texttt{waypoint\_sender} and \texttt{battery\_monitor} nodes.

Each of the four readiness gates described above follows the same general pattern of an unbounded polling loop with a fixed wait interval between attempts. The implementation avoids a timeout condition entirely, and the only failure condition for a readiness gate is the failure of its subprocess. This approach deliberately avoids adding an additional parameter with a timeout for each readiness gate. At the moment when such a timeout would be reached, the launch tree would need to make a decision about what to do next: fail the entire launch, proceed without this supposedly failed prerequisite, or retry the check later~\cite{blass2021automatic}. All of these options, however, reintroduce the same class of assumptions about timing that the original \texttt{slam.launch.py} sought to eliminate by using \texttt{ros\_gz\_bridge} and hardcoded delays in the first place. Instead, readiness gates poll without bound and, in the presence of a genuine failure, simply cause the launch tree to linger indefinitely at a known location~\cite{lampe2023robotkube}. This behavior is desirable from a debugging perspective, as the poller subprocess output to stdout is now the only source of information about what went wrong. The same log line can be used to determine which readiness condition is not met, and thus the cause of the launch tree hanging. A potentially unbounded delay before mission startup is exchanged for the ability to diagnose intermittent issues of this class, which is broadly consistent with the reliability philosophy articulated in Chapter~6~\cite{lampe2023robotkube}. The complete implementation of these deterministic polling primitives is provided in Appendix~\ref{app:code_launch} (Listing~\ref{lst:gate_primitives_core}).

\subsection{OnProcessExit Daisy-Chaining and the Per-Robot Dependency Graph}

Readiness gates alone do not form a launch architecture; they have to be wired into the launch framework's action graph so that a gate's exit triggers the next set of processes that it was supposed to enable. \texttt{multi\_robot\_controlled\_launch.py} uses \texttt{RegisterEventHandler} bound to \texttt{OnProcessExit} for that purpose - a launch-framework primitive that fires an event handler whenever the target process exits due to any reason. At the same time, since all the gate processes in this design exit cleanly with exit code 0 only when their readiness condition is met, and hang otherwise to keep the launch sequence from proceeding, attaching an \texttt{OnProcessExit} handler to a gate is effectively equivalent to defining an action that should take place when this gate is done, for as long as the launch framework knows about such a thing as a readiness condition.

The global bring-up sequence is then constructed in a linear fashion around the world gate: The clock bridge, clock gate, and mission-gate node all start right after the world gate is closed, with the individual robot chains then diverging from robot1's spawn process. \texttt{\_robot\_chain} is a factory function that returns the list of launch actions specific to a robot, plus a reference to the corresponding lifecycle gate that serves as the trigger for whichever robot chain is supposed to follow it. Within a single robot's launch chain, there are then four sequential links: The clock gate (or the previous robot's lifecycle gate, for robot2 and robot3) trigger spawn, the spawn process triggers the sensor bridge, robot-state publisher, EKF filter node, and topics gate, the topics gate in turn triggers the full Nav2 stack plus the lifecycle gate, and finally the lifecycle gate triggers the two mission-level nodes, \texttt{waypoint\_sender} and \texttt{battery\_monitor}.

The reason for deferring robot2 and robot3 to subsequent clock-gate dependencies rather than simply launching all three chains in parallel from the clock gate is not a concurrency control measure - at least not directly. As far as the logical launch sequence is concerned, robot2's chain could begin before robot1 reached an active Nav2 lifecycle state - but it does not, not because of any robot-specific scheduling policy, but rather due to limitations in the WSL2 host that runs the simulation. There is only so much CPU time to go around, and three simultaneous gazebo spawns, three simultaneous nav2 stacks each with six lifecycle nodes, and three separate EKF filters each fusing odometry and IMU information concurrently, taxes the CPU to the point that action servers to timeout erroneously in the middle of bring up rather than cleanly and predictably failing. The alternative of launching all three robots in parallel simply does not work due to resource constraints, which is why the three robot chains are launched one after another: robot2's spawn is deferred until robot1 reaches an active lifecycle state, and robot3's spawn is deferred until robot2 reaches that same state. This sacrifices total bring-up time to keep instantaneous CPU consumption within limits during the bring-up, and the same design consideration likely underlies the choice to launch the robots' chains sequentially in general, rather than interleaving the lifecycle gates for different robots. It is the analog, within the confines of a single WSL2 host, of the scheduling policy choice between sequential and parallel execution described more generally in the distributed-systems literature on task scheduling under bounded shared resources.~\cite{blass2021ros}.

Another consequence of launching everything via \texttt{OnProcessExit} events from the launch framework's action graph, rather than from a central scheduler with state about the robots being launched, is that the dependency graph does not have any information about which robot a certain step refers to, beyond the closures captured when the steps were defined. Each call to \texttt{\_robot\_chain} binds the factory-closed variables to specific values: the bridge YAML, the Nav2 parameters, and the spawn position for the robot that is being launched, while the action that is supposed to trigger this particular set of actions is simply whatever gate or lifecycle-gate binding the caller chooses to pass to it. The graph is effectively built in an ad-hoc procedural manner inside the \texttt{OpaqueFunction}-wrapped \texttt{\_build\_chains}, because the name of the Gazebo world, and consequently the world gate, is not known to the launch framework until the world file path is itself parsed at the point the launch context is evaluated.

One design choice that stands out as particularly thoughtful is the decision to bind the RViz visualization node to robot3's lifecycle gate rather than to the clock gate. By making RViz start only after the third robot has entered an active lifecycle state, the visualization tool does not attempt to display costmaps and TF transforms for robots that have not yet spawned - reducing visual clutter during bring-up before these topics are published. The gate-based deterministic launch orchestration architecture, utilizing custom readiness polling primitives and \texttt{OnProcessExit} event handlers to sequence multi-robot bring-up without race conditions, is shown in Figure~\ref{fig:ch7_launch_architecture}. The programmatic construction of this launch dependency graph and its sequential per-robot lifecycle daisy-chaining is specified in Appendix~\ref{app:code_launch} (Listing~\ref{lst:launch_chains_core}).

\begin{figure}[htbp]
\centering
\includegraphics[width=\textwidth]{fig_7_1_launch_architecture.pdf}
\caption{Deterministic gate-based launch architecture depicting Gazebo world loading, clock bridging, topics readiness verification, and sequential per-robot lifecycle daisy-chaining.}
\label{fig:ch7_launch_architecture}
\end{figure}

\subsection{Idempotent Spawning and Stale Process Cleanup}

Two additional robustness features address the possibility of repeated invocations of the same launch graph - for example, when development iterations require multiple bring-ups without the intermediate step of shutting down the entire development environment. \texttt{spawn.py} provides an idempotency guard around the actual \texttt{ros2 run ros\_gz\_sim create} command, while \texttt{slam.launch.py} performs a global sweep for stale Cartographer processes before it attempts to launch a new one.

The need for the idempotency guard in \texttt{spawn\_robot} comes from the fact that a launch tree constructed from \texttt{OnProcessExit} handlers, as per Section 7.2.3, has no awareness of previous invocations of the same command. In particular, if a developer forgets to shut down the gz server process before re-invoking the launch file, an unprotected \texttt{ros\_gz\_sim create} call would either fail with a name collision if the robot's model name is not unique across runs, or, if it is unique, silently launch a second robot alongside the first. Neither outcome is desirable, and not simply because of the wasted CPU resources: A development iteration that requires a fresh bring-up to be performed multiple times is more likely than not to involve the entire suite of launch files being re-invoked, including ones that make no attempt to clean up after themselves.

Beyond the idempotency check, the spawn function implements an unbounded retry loop for the create command, rather than a single attempt followed by failure, and additionally polls \texttt{gz model --list} independently after a successful create return code, rather than relying on that return code to indicate success. The reasoning is that \texttt{ros2 run ros\_gz\_sim create} returning exit code zero only indicates that the request to spawn has been accepted for execution by Gazebo, but not that the request has completed, i.e., that the physics engine has instantiated the model and made it queryable. Attempting to treat the exit code of the create command as indicative of spawn completion would re-introduce the same discovery-versus-delivery gap as the clock gate blocking echo described in Section 7.2.2: a request to perform an action has been accepted, but the effects of that action have not yet been realized, and subsequent launch stages should therefore wait until the action is known to be complete.

A similar consideration applies to the cleanup of previously orphaned Cartographer processes performed before launching the new ones. Since Cartographer's nodes are launched as plain old \texttt{ExecuteProcess} nodes, not as lifecycle nodes, sending SIGINT to the main ROS 2 launch process does not kill the child Cartographer processes reliably - in particular, under WSL2 it often leaves behind two processes bound to the same \texttt{/map} topic. A subsequent attempt to launch Cartographer would then result in two processes fighting for control of the \texttt{/map} topic, with subsequent effects ranging from intermittent map corruption in RViz to outright launch-file crashes that need SIGINT to clear. However, the cause of such a failure is not obvious at the time - if the first invocation of \texttt{slam.launch.py} was interrupted, the second one's failure to initialize may wrongly be attributed to a bug in the launch file itself, when in reality the problem was the leftover processes from the first invocation.

The cleanup is implemented in the form of an \texttt{OpaqueFunction}, placed at the beginning of the launch description, that runs\texttt{pkill -f} to kill any processes matching the regular expression patterns of the full path to \texttt{cartographer\_node} and \texttt{cartographer\_occupancy\_grid\_node}. This approach to the stale-process cleanup exemplifies a recurring theme in this particular launch file - that of making it possible to perform common maintenance without the need for extensive process discovery. If the same code were to be used to check if a stale process exists before attempting to kill it, the same \texttt{pkill -f} construct would suffice - it simply kills all existing processes matching the pattern, but if there are none, then there is nothing to do. Meanwhile, an alternative process-removal approach based on checking first would require inspecting \texttt{/proc} to perform the same maintenance, but for no particular reason other than to report that no process was found.~\cite{shuai2026roscell}.

\section{Namespace Isolation and Multi-Robot TF Tree Construction}

\subsection{URDF Link Prefixing via Regular-Expression Patching}

To configure three TurtleBot3 Waffle units within a shared TF tree, each link and joint frame needs to use robot-specific names, since having three \texttt{base\_link} frames is incompatible with the tf2 tree topology. One apparent solution is to use the \texttt{frame\_prefix} parameter for \texttt{robot\_state\_publisher}, which prepends a given string to every frame name published by the node. \texttt{\_patch\_urdf} in \texttt{multi\_robot\_controlled\_\-launch.py} deliberately avoids using the parameter, instead applying necessary transformations to the URDF source string itself before passing it to \texttt{robot\_state\_publisher}. The comment in the module documents the decision: using \texttt{frame\_prefix} would "double the prefix," making manual patching preferable.

The issue arises because \texttt{robot\_state\_publisher} does frame prefixing for every frame name parsed from the URDF, including those referenced implicitly by the contents of \texttt{<parent>} and \texttt{<child>} tags; if the prefix is added a second time by another component also parsing the URDF, having \texttt{robot1/base\_link} appear in some parts of the code and \texttt{robot1\_robot1/base\_link} appear in others becomes a possibility. Instead of trying to reason about which frame-name construction method takes precedence under different circumstances, patching the URDF source directly makes any frame name using the robot-specific prefix consistent across all usages~\cite{kaiser2022ros2swarm}. This avoids subtle bugs downstream where a frame name constructed using string concatenation elsewhere in the code has a different value than the one published by \texttt{robot\_state\_publisher}.

This makes \texttt{\_patch\_urdf} apply the same transformation to every occurrence of \texttt{<link name="...">}, \texttt{<parent link="...">}, and \texttt{<child link="...">} tags, adding the robot-specific component to the parsed link name and leaving the syntax of the tag itself intact. The function implements a guard to avoid adding the prefix twice when testing the patched URDF string, should an accidental second application of the patch occur; this is a defensive programming technique to avoid subtle bugs elsewhere, as failing to skip an unintended double-application of the patch would manifest much later as an incorrect TF tree structure.

The patching function itself only targets the three patterns listed, instead of attempting to rename all attributes or tags that reference a link name with a regular expression. This is because, in the particular URDF describing the TurtleBot3 Waffle used in this project, visual, collision, and inertial tags only appear inside a \texttt{<link>} definition and do not need renaming, and no other tags reference a link name by string outside the three patterns. A general-purpose XML parser is required to implement such a rename for an arbitrary URDF file, since valid URDF may also make use of Xacro macros or \texttt{<gazebo>} tags containing additional link references; for the particular known robot description used for the TurtleBot3 Waffles, however, a regex-based solution is sufficient to implement the necessary transformations without the overhead of an additional parser~\cite{kaiser2022ros2swarm}.

\subsection{Double TF Publisher Conflict Elimination}

The multi-robot TF tree management also encounters a second issue related to namespace isolation: the frame \texttt{odom} $\rightarrow$ \texttt{base\_footprint} has two potential publishers in the default Gazebo differential drive configuration, the simulated diff-drive plugin and the \texttt{ekf\_filter\_node} sensor fusion filter. \texttt{bridge\_robot1.txt} comments section documents the problem before demonstrating the solution: "the bridge block for relaying the plugin's own TF was completely removed to eliminate double TF publisher conflicts." The EKF node itself becomes the sole publisher for the transform.

Having two nodes attempt to publish the same parent-child TF frame relationship is not a situation Nav2 is prepared to handle: since tf2 does not support resolving writer conflicts, the two transforms simply end up being published in parallel, with any downstream consumers having to deal with whichever transform message arrives first. In the particular case of the combination of the diff-drive plugin and the EKF sensor fusion filter, having the EKF recalibrate wheel-odometry covariance as recommended in Chapter~5, Section~5.4 causes the filtered transform to alternate unpredictably between the two independently-published pose estimates as each publisher's message happens to arrive; when the two transforms differ, a downstream component such as \texttt{amcl}'s motion model updating step will encounter a discontinuity in received odometry, causing momentary degradation of particle cloud quality unrelated to sensor noise properties and unsuspected by AMCL.

This creates a situation similar to the "single-writer" rule for shared, concurrently-accessed variables in the context of concurrency control; instead of two writers for a shared variable resulting in undefined behavior, having two nodes publish the same transform results in intermittent transform failures downstream. This can be resolved in the same manner as such a scenario would be resolved for shared memory: by eliminating one of the writers, so that a single transform publisher remains. This is what the commented-out bridge section for the diff-drive plugin accomplishes, removing the double publisher conflict by disabling the plugin's TF broadcaster entirely. The same concept is used in the "spatial-claim" section of the mutual exclusion algorithm described in Chapter~6, used for resolving which of two robots attempting to enter a warehouse corridor simultaneously has claimed it first; by designating the corridor segment as claimed by default and only allowing a second robot to enter if the first has released the claim, the potential for simultaneous entry by both is eliminated~\cite{abaza2025aicovariance}.

The single-writer policy for TF transform publishers has one particular implication for the launch configuration: just like in the mutual exclusion algorithm, the writer that holds the exclusive claim for the transform must be up and publishing before any downstream components, such as \texttt{amcl} or Nav2's costmap, attempt to consume it. This becomes the responsibility of the topics gate described in Section~7.2.2, since the topics readiness list for each robot definition contains the filtered odometry topic published by the EKF node along with the scan and joint\_state topics. By virtue of using that particular readiness list, the launch configuration ensures that Nav2's activation will only take place after the EKF node has started publishing the transform.

\subsection{Global vs. Namespaced Topic Boundaries and the DDS Isolation Workaround}

The per-robot Nav2 stack, defined in Section~7.3.2, runs within a \texttt{GroupAction} that is wrapped in a \texttt{PushRosNamespace(name)}, which ensures that any topics, services, or actions created or discovered by relative names within a robot's nodes are namespaced under /robot1, /robot2, or /robot3. This is the correct and necessary behavior for the vast majority of topics these robots will publish to or subscribe from: three instances of \texttt{amcl} publishing to a relative \texttt{amcl\_pose} topic must appear on three different namespaced topics or the whole point of running these stacks concurrently is defeated. The fleet-coordination layer built upon this namespaced foundation, however, requires the converse behavior for a small minority of topics, and the code comments reveal where this design flaw was recognized and addressed.

\texttt{TaskCoordinator.add\_waypoints\_pub} in \texttt{waypoint\_sender.py} is commented to "Publisher changed to absolute/global to resolve DDS namespace isolation bug," with the subscription in the same file marked with a similar observation. The reason for this change is that \texttt{/add\_waypoints} and \texttt{/side\_claims} are not per-robot topics at all: the battery-proportional task-reallocation message of Chapter~6, Section~6.6 and the corridor-claiming topic of Section~6.5 must be broadcast to all surviving robots, regardless of which robot initially published them. Declaring these publishers and subscriptions to exist within the same relative-name convention as the rest of the robot's interfaces would cause the \texttt{PushRosNamespace} wrapping the \texttt{TaskCoordinator} \texttt{GroupAction} to rewrite them to something like \texttt{/robot1/add\_waypoints}, \texttt{/robot2/add\_waypoints}, and \texttt{/robot3/add\_waypoints}, resulting in three topics with the same name suffix but no cross-publication between them, since no robot's node is namespaced to listen on any namespace but its own. The DDS-level topic-discovery mechanism would not report this as an error at all: within each robot's own namespace, that robot's publisher and that same robot's own subscriber would match each other perfectly, appearing to DDS as a fully healthy, correctly discovered topic pair. What never occurs is any match \emph{across} the three per-robot topics, since the entire purpose of the message is to reach robots other than the one that published it. This design error would manifest itself by task-reallocation messages being published into a namespace no other robot subscribed to, with no error message at all being generated by the middleware.

The workaround used in the code comments, declaring \texttt{/add\_waypoints} and \texttt{/side\_claims} with an absolute, namespace-free path, instructs ROS~2 to treat these topics as already fully qualified and prevent any parent \texttt{PushRosNamespace} from modifying them. This resolves the isolation bug while minimally modifying the per-robot interface design of Section~7.3.1 and Section~7.3.2: the same \texttt{GroupAction} wrappers are used throughout, ensuring relative-name topics follow the namespace isolation convention while \texttt{/add\_waypoints} and \texttt{/side\_claims} explicitly escape it. The isolation bug workaround is minimally invasive in this sense and only affects the topics that require it, becoming another example of the ROS~2 design principle of wide applicability for small changes~\cite{macenski2022robot}. It is noteworthy from a design perspective as well, as it avoids the alternatives of revising the namespace layout strategy across all robots or writing a dedicated relay node to forward messages between \texttt{/robot1/add\_waypoints} and \texttt{/robot2/add\_waypoints}, or indeed any combination of the three.

\section{Diagnostic Instrumentation and Passive Telemetry Architecture}

\subsection{Dual-QoS Subscription Contracts for Passive Observability}

Chapter~5, Section~5.2 established in the abstract that the diagnostic framework maintains "both volatile best-effort $\mathcal{Q}_{\text{BE}}$ and transient-local reliable $\mathcal{Q}_{\text{TL}}$ contracts" against AMCL pose publications, while the active inference agent subscribes exclusively to the reliable transient-local profile. \texttt{amcl\_diagnose.py} is the concrete realization of that theoretical contract: it defines \texttt{AMCL\_QOS}(depth=10, VOLATILE, BEST\-\_EFFORT) and \texttt{AMCL\_QOS\_ALT}(depth=10, TRANSIENT\_LOCAL, RELIABLE), subscribes to \texttt{/robotN/\-amcl\_pose} twice per robot with the respective QoS contracts, and maintains both subscriptions concurrently.

The necessity to maintain two subscriptions to the same topic stems from the ROS2 DDS middleman's behavior when a subscriber specifies a QoS stricter than a publisher's: instead of rejecting the subscription attempt, the subscriber silently fails to receive messages from that publisher. \texttt{nav2\_amcl} advertises its \texttt{amcl\_pose} topic with a particular (but possibly varying between deployments) contract; if the subscriber's contract differs from the publisher's, the subscriber receives no messages. Since the entire raison d'etre of a passive observer is to obtain trustworthy information about the observed system even when the observed system cannot be relied upon to provide it, a diagnostic subscriber that fails to receive AMCL messages due to a contract mismatch silently fails its primary mission. Subscribing to the same topic with two different QoS profiles ensures that the diagnostic subscriber receives AMCL messages regardless of which convention the amcl publisher chooses, but incurs the computational cost of having AMCL messages processed twice on either subscriber callback: \texttt{SampleCollector.\_amcl\_cb\_factory} does not handle that scenario as an error to be guarded against, but rather accounts for it as a possibility to defend against by checking the received timestamp against \texttt{\_last\_amcl\_stamp} for that robot and deduplicating samples whose timestamps are no earlier than $10^{-9}\,\text{s}$ difference, well within the range of physically meaningful AMCL update frequencies and safely beyond floating-point precision errors. This approach to handling the subscription to a potentially-unpredictable publisher under two contracts acknowledges the increased robustness of the "subscriber knows better" approach as a valid design goal, but rather than denying the possibility that the subscriber might be wrong, it avoids the cost of guarding against it by instead making it a localized consideration that does not affect the diagnostic logic's integrity.

The same dual-subscription technique is used to address a separate concern unique to the preflight diagnostics: \texttt{amcl\_rx\_count} gets incremented every time the subscriber receives a message on either of the two subscriptions, but before attempting to extract relevant information from the message, the collector uses the same \texttt{amcl\_rx\_count} to determine the number of messages that have been passed to the diagnostic record, and \texttt{amcl\_parse\_errors} and \texttt{amcl\_parse\_first\_error} to track malformed messages (those with an incomplete covariance array or an unparsable timestamp). Thus, the diagnostic report tracks three independent metrics related to message reception: messages received (rx\_count), messages processed (stored in the record), and messages parsed successfully (implied by the lack of increment to amcl\_parse\_errors). If a particular robot's report shows unexpectedly high rx\_count but low or zero entries in any of the stored pose records, the discrepancy suggests that either an upstream component was generating extraneous AMCL messages, or the diagnostic subscriber was receiving those messages concurrently over two subscriptions. If the rx\_count is high, with a nonzero amcl\_parse\_errors count, then the problem is specific to individual malformed messages that the diagnostic tool was able to reject on parsing attempts but not on initial subscription. The distinction between the three concepts allows the diagnostic report to highlight three independent failure modes instead of one binary outcome of either successful message reception or not~\cite{lee2025dependency}.

A similar observation can be made regarding the \texttt{preflight\_probe\_amcl} component that determines whether any AMCL traffic whatsoever is visible to the subscriber before a live collection is initiated. The probe uses the same dual-subscription approach to ensure visibility into AMCL activity regardless of the publisher's contract, but collects the statistics over a fixed $2.0\,\text{s}$ interval before terminating, its own docstring framing the function as a quick probe of whether AMCL messages are visible to the process at all rather than as a characterization of their rate. The preflight collection's sole utility is to report whether the upcoming live collection will fail immediately upon launch due to an invisible AMCL publisher; the probe's statistics only need to be indicative of the situation at the moment of the probe's execution, and its fixed $2.0\,\text{s}$ interval can report a failure if the observed count across two subscriptions is nonzero, regardless of that count's magnitude. The distinction is critical: a sufficiently long preflight collection window would report more traffic than useful to the diagnostic tool, but the \texttt{preflight\_probe\_amcl}'s purpose is to not analyze long-duration AMCL statistics, but rather to confirm that an AMCL publisher that the diagnostic tool subscribes to is visible to it at all, which it accomplishes by collecting statistics at a fixed rate for a fixed interval right before launching the actual diagnostic collection, rather than attempting to interrupt that collection to assess the subscription status mid-way through.

\subsection{Event-Driven Flight Recording vs. Fixed-Rate Telemetry Sampling}

\texttt{diagnosis.\allowbreak py} and \texttt{localization\_\allowbreak diagnose.py} collect telemetry about the same physical quantities (\allowbreak AMCL poses and covariances) but use fundamentally different collection approaches: the diagnosis module uses an event-driven state tracking, only issuing reports when a robot's state changes, while \texttt{localization\_\allowbreak diagnose.py} uses a fixed-rate sampling to build dense telemetry records about the fleet's AMCL stats over time. That difference in approach is codified in the \texttt{diagnosis.py} module docstring, which labels the current version "v1.1 - Event Driven" and explicitly states that this iteration "eliminates 10Hz logging noise," implying a superseded fixed-rate predecessor without naming it directly.

\texttt{MultiRobotSystemDiagnosis} keeps the $1\,\text{Hz}$ \texttt{diagnostic\_timer}, but all it does is drive a silent boundary and proximity check inside the \texttt{\_diagnostic\_tick} callback: all of the actual report output happens in the \texttt{\_robot\_status\_callback} handler, which only calls \texttt{\_print\_system\_state\_summary} when the robot's status has changed compared to the previous state; the claim, completion, and deferral handlers follow the same pattern, only printing a state summary when a change occurs. As a result, the only content in the log output for a robot that spends the majority of its mission time in an idle state is the initial log entry, followed by entries whenever that robot transitions out of (or back to) idle and the corresponding time spent since the last update. For robots that spend time performing active localization, the report is denser, reflecting the frequency of updates to their pose estimate, but for those that spend the majority of their runtime in passive localization, the report is sparse to match.

\texttt{LocalizationDiagnosticSuite} takes the opposite approach: \texttt{\_record\_telemetry\_tick} is driven by a $0.5\,\text{s}$ fixed-period timer, explicitly noted in the code to "avoid CPU overhead on WSL2" compared to the $10\,\text{Hz}$ rate used by \texttt{diagnosis.py}, and on every tick, it appends a new record to the robot's telemetry arrays regardless of whether any of the recorded values have changed since the last update. The report only contains the timestamps, TF latency, yaw offset, and covariance sigma for every recorded tick, but the record's density is required for the final report's statistic generation: the mean and maximum for each of the recorded values are only meaningful as descriptions of a uniform distribution if the underlying data consists of measurements taken at regular intervals across the mission's runtime.

The different collection approaches in \texttt{diagnosis.py} and \texttt{localization\_diagnose.py} are not contradictions but rather consequences of the type of report each one generates. The former's objective is to report a discrete mission timeline of robot claims, waypoints, and deferrals as a series of events in chronological order, with details only provided for the actual mission state changes, while the latter aims to report continuous statistics about the fleet's behavior based on samples taken at regular intervals throughout the mission. It would be wasteful to populate the event-driven report with entries for every fixed-period tick, and equally uninformative to attempt to generate statistics about a uniform distribution from the event-driven report's sparse records, but both approaches are entirely justifiable within their context. Recording timestamps for every discrete change in the mission state provides an informative temporal context for each event in the mission timeline, while fixed-period sampling provides the density necessary for calculating mean and maximum values for distributions such as TF latency and AMCL covariance drift. Attempting to answer the question of "what happened during the mission" using only timestamps from the event-driven approach would result in a report with long stretches of silence between consecutive events, while trying to populate the event-driven report with timestamps on a fixed schedule would recreate the "10Hz logging noise" that \texttt{diagnosis.py} was explicitly changed to eliminate~\cite{blass2021automatic}.

\subsection{Causal Correlation of Debug Events against Nearest-Sample Estimator State}

The diagnosis package's \texttt{amcl\_diagnose.py} script implements a second, offline capability that parses existing debug logs for discrete textual events occurring in its per-robot debug logs (e.g., \texttt{SENT\_GOAL}, \texttt{goal rejected}, \texttt{collision ahead}) and, for each, determines what the AMCL, odometry, and TF state of the corresponding robot most likely was at the time. Since the debug log and AMCL sample record are two largely independent data sources cross-referenced only by a shared, inconsistently sampled time reference axis, this requires an explicit \texttt{nearest\_before} correlation rule: for any given debug-log event time, find the last AMCL sample with a timestamp not later than it and ignore any later samples, since those occurred after the event and thus cannot have caused it. This is an intentionally asymmetric choice: had the rule been to find the nearest sample in either direction, an event with a timestamp in between two samples would have had two possible matches, one before and one after, with the later one being recorded as the one that occurred first, thereby defeating the diagnostic value of distinguishing causes from effects.

Only the most recent preceding sample is selected to allow for any correlated row to have a cause-effect relationship strictly in one direction (per the physics of the system under test), even at the cost of occasionally having to correlate an event with a several-seconds-old sample if the AMCL publish rate happened to be particularly slow at the time. This potential several-seconds staleness is also explicitly exposed as part of \texttt{amcl\_age}, the signed difference between the event timestamp and the matched sample timestamp, and flagged as \texttt{bad\_age} in \texttt{summarize\_rows} whenever this staleness exceeds $2.0\,\text{s}$; similarly, \texttt{bad\_sigma} is flagged if the matched sample's sigma in position is worse than $0.35\,\text{m}$. This terminology is deliberately consistent with the thresholding semantics of the live AMCL recovery mechanism from Chapter~6, since \texttt{amcl\_diagnose.py}'s offline report is effectively asking the same question that mechanism answers online: had $\sigma_{\text{amcl}}$ exceeded threshold at this moment during an actual mission run, would \texttt{waypoint\_sender}'s \texttt{\_trigger\_amcl\_recovery\_spin} have fired? For the live system, this is decided directly by the threshold check inside \texttt{\_amcl\_callback}; for \texttt{amcl\_diagnose.py}, the same question is answered retrospectively by correlating a debug-log event with a stale AMCL sample and reporting it as \texttt{bad\_sigma}.

A subtle correctness concern specific to this script is that, in the debug logs being parsed for events, timestamps might have been recorded using a different clock convention than the one AMCL is using to record the samples being parsed in a later run. A raw epoch timestamp recorded in a debug log from a previous development session and a simulation timestamp recorded in the bag from the current one can easily differ by several orders of magnitude, depending on how the clock was configured at the time, not because of any localization staleness but simply because the clocks were measuring different things entirely.

\texttt{correlate\_events} handles this by treating such clock domain mismatches as invalid and \seqsplit{non-localization-related}, and, if the raw age computation would have failed to satisfy $-120.0\,\text{s} < \text{age} < 120.0\,\text{s}$, as a threshold chosen to be comfortably larger than the maximum plausible publish interval for any AMCL publisher but comfortably smaller than the clock domain mismatch, skipping the invalid row and marking \texttt{clock\_domain\_mismatch} instead of either \texttt{bad\_age} or \texttt{bad\_sigma}. This guards against trying to interpret the clock domain mismatch as a localization staleness by erroneously reporting an age of several hours, which \texttt{summarize\_rows}'s hard-coded threshold of $2.0\,\text{s}$ would then correctly flag as \texttt{bad\_age} but for the wrong reason. This separation of concerns - having localization staleness checks and clock domain mismatch checks implemented as separate, orthogonal rules rather than conflating the two into a single large timestamp age check - is particularly useful in diagnosing why a particular localization staleness finding was reported, as it prevents the diagnostic report from conflating two separate issues in the process. In the broader context of the localization diagnostics, however, it means that reusing historical debug logs generated by earlier versions of the software across development sessions can generate misleading clock\_domain\_mismatch reports~\cite{kudla2026system}.

\subsection{Bag-Based Post-Hoc Consistency Verification}

While \texttt{amcl\_diagnose.py} focuses on correlating individual debug events against the AMCL estimator state at the nearest preceding sample, the \texttt{analyze\_bag.py} script is concerned with an equally important but subtly different question: given a bag recording of the full set of odometry, AMCL, and TF messages for a completed run, is the filtered odometry trajectory consistent with the estimated global pose over the recorded time interval? This transforms the diagnostic focus from looking at instantaneous localization estimator states in response to individual debug events into examining the consistency of the estimator's trajectory as a whole, and the two scripts differ starkly in both their design and implementation details: \texttt{amcl\_diagnose.py} parses a debug log for events and correlates them against estimator samples one by one, while \texttt{analyze\_bag.py} has to process three independently recorded but time-ordered streams of messages with timestamps in \texttt{rosbag2\_py}.

The inner loop consists of composing the odometry-frame position with the \texttt{map} $\rightarrow$ \texttt{odom} correction transform recovered from \texttt{/tf} at the same timestamp and comparing the result to the pose published on \texttt{amcl\_pose} using the same nearest\_before binary search as \texttt{amcl\_diagnose.py}. For an odometry-frame position $(x_{\text{odom}}, y_{\text{odom}})$ and a map-to-odom correction $(map_x, map_y, map_{\text{yaw}})$ recovered from the nearest preceding \texttt{/tf} message, the predicted global position is given by
\begin{equation}
\begin{aligned}
x_{\text{pred}} &= map_x + x_{\text{odom}} \cos(map_{\text{yaw}}) - y_{\text{odom}} \sin(map_{\text{yaw}}), \\
y_{\text{pred}} &= map_y + x_{\text{odom}} \sin(map_{\text{yaw}}) + y_{\text{odom}} \cos(map_{\text{yaw}})
\end{aligned}
\label{eq:predicted_global_position}
\end{equation}
and the discrepancy, $\sqrt{(x_{\text{pred}} - x_{\text{amcl}})^2 + (y_{\text{pred}} - y_{\text{amcl}})^2}$, between this position and the simultaneously published \texttt{amcl\_pose} is logged for all messages and summarized by their mean and 95th percentile in the final summary.

A large discrepancy at this stage can be indicative of either of two problems: that the \texttt{map} $\rightarrow$ \texttt{odom} correction published to \texttt{/tf} and the \texttt{amcl\_pose} have diverged within AMCL's own publication logic, or that the odometry frame the correction is being composed against has itself diverged substantially from the corrected value. Either of these is grounds for concern, but distinguishing between the two requires further analysis beyond the scope of this script. At a higher level, this summary discrepancy statistic represents a coarse, single-number sanity check for the consistency of the estimator's published pose with the odometry frame it is using to publish it, and such a check provides a useful precursor to any more detailed analysis of why the discrepancy might be occurring. The finer-grained event-based analysis performed by \texttt{amcl\_diagnose.py} and described in detail in Section~7.4.3 is, in fact, a natural follow-up to this bag-level sanity check, as it attempts to diagnose the discrepancy by looking at whether any of the events explicitly logged by the node under test as of significance could plausibly be responsible for it; likewise, a more detailed analysis of the odometry drift itself would benefit from the tools and insights described in Section~7.4.1.

Finally, the \texttt{bag\_to\_csv.py} script subscribes to \texttt{/robot1/odom}, \texttt{/robot1/amcl\_pose}, and \texttt{/gazebo/\-model\_states}, writing each of the three streams to a separate CSV file for later manual inspection. Unlike the other two scripts, it does not attempt to perform any correlation or analysis of the three streams against each other, instead only serving to extract these message streams, which are not otherwise directly accessible, into a more accessible format for ad hoc analysis or plotting. This functionality is particularly useful for the \texttt{/gazebo/model\_states} stream as a source of ground truth for the robot pose, which the \texttt{analyze\_bag.py} script does not make direct use of but which may, in some cases, be useful to compare against both the odometry and the estimated \texttt{amcl\_pose}. Storing the raw CSV data separately from the analysis results, as this script does compared to, say, having \texttt{bag\_to\_csv.py}'s CSV-writing logic folded directly into \texttt{analyze\_bag.py}, serves to decouple the end user's direct inspection of all the data from whatever higher-level statistical analysis the end user may want to perform separately on a subset of those data~\cite{kudriashov2026real}.

\section{Calibration and Auxiliary Tooling}

\subsection{IMU Frame Relay and Odometry Covariance Recalibration in Deployment}

Chapter~5, Section~5.4 calculates the calibrated odometry covariance matrix $\boldsymbol{\Sigma}_{\text{calibrated}}$ of Equation~\eqref{eq:calibrated_covariance_matrix} as a corrective measure to the overconfident, near-zero variance default values the Gazebo diff-drive plugin publishes by default. This calculation in turn explained why the calibrated values take the particular magnitudes they do; this section instead deals with the narrower question of where in the system this correction ultimately gets applied, and why this correction gets implemented as a standalone tool rather than a parameter inside the diff-drive plugin itself.

In the deployed setup, \texttt{odom\_cov\_adjuster.py} subscribes to a robot's \texttt{/{robot}/odom} topic and republishes an otherwise identical Odometry message to \texttt{/{robot}/odom\_fixed}, with the \texttt{pose.covariance} field overwritten with the \texttt{DEFAULT\_COV} array from $\boldsymbol{\Sigma}_{\text{calibrated}}$. It is worth noting that the currently deployed configuration does not make use of this relay; the per-robot Nav2 parameter file does not include any parameters for \texttt{ekf\_filter\_node} related to scaling the covariance, while the node's parameters only concern its update rate, frame, and the two topics it subscribes to, along with a per-axis mask determining which state axes get fused from which estimator. Instead, the overconfidence the calibrated covariance of Section 5.4 sought to resolve gets mitigated by another change specific to the robot's estimator configuration: the mask for texttt{/robot1/odom} only trusts the linear velocity along the forward axis, and the angular velocity around the vertical axis, discarding the position and orientation fields~\cite{tripathi2023sensor} the wheel encoders have little capacity to estimate reliably over distance. The covariance scaling and axis masking represent two sides of the same problem, with the deployed configuration only utilizing the latter, while \texttt{odom\_cov\_adjuster.py} remains an unregistered utility script implementing the former, a discrepancy this report notes but does not seek to resolve, since it would entail either registering the relay as a console-script node or adding covariance-scaling parameters to the current configuration.~\cite{tripathi2023sensor}.

This discrepancy represents less of a conflict and more of a case of retaining a debugging utility from earlier development stages as a standalone tool. \texttt{odom\_cov\_adjuster.py} remains useful in its current capacity as a way to recalibrate and observe the downstream effects of a different EKF covariance in isolation, without any changes to whichever node consumes the \texttt{/{robot}/odom\_fixed} topic, be it an EKF or another tool. A standalone relay represents the correct way to isolate and test a hypothesis, be it recalibrating the covariance as Section 5.4 does to determine its downstream effects before committing to a change in the configuration, or retaining such a tool as a utility script to observe those effects in a deployed setup~\cite{abaza2025aicovariance}.

A related but stand-alone relay, \texttt{imu\_relay.py}, executes the same kind of republishing of the IMU stream, subscribes to \texttt{/imu} and republishes to \texttt{/imu\_fixed} with the \texttt{header.frame\_id} overwritten to \texttt{imu\_link}. Unlike the case of the odometry relay, this change is not part of the per-robot mission chain at all: \texttt{multi\_robot\_controlled\_launch.py}'s \texttt{\_robot\_chain} launches the bridge, robot-state publisher and the EKF node directly, without any \texttt{imu\_relay} participant, and the \texttt{ekf\_filter\_node} itself subscribes to the raw, unrebatched \texttt{/robot1/imu} topic as its \texttt{imu0} input. The relay's actual consumer is the single-robot Cartographer mapping pass of Section 7.7.2: \texttt{slam.launch.py} explicitly launches \texttt{imu\_relay} and remaps \texttt{imu} to \texttt{/imu\_fixed} before launching \texttt{cartographer\_node}, since the latter's IMU-fused trajectory builder, as detailed in Chapter 5, Section 5.2 when discussing the SLAM parameter set, requires a properly relabeled IMU frame to perform angular velocity fusion in its scan matcher.

Both relays thus effect tweaks to a default-topic publication, either by modifying frame labels or adjusting covariances, and both are implemented as standalone scripts rather than built into the bridge configuration, which could have been an option for either one given their current status as unifying single- and multi-robot setups in this workspace. The reason for this choice is that, in contrast to the odometry-covariance scaling question, the IMU relay's \texttt{/imu\_fixed} tweak is already fully integrated into the separate mapping pipeline it was developed alongside.~\cite{moore2016generalized}.

\subsection{Interactive Yaw Calibration for SLAM Map Frame Alignment}

\texttt{custom\_slam\_teleop.py} defines a teleoperation loop that only gets used once, during the Cartographer offline mapping pass of Section~7.7.2, and serves a singular calibration-related purpose: aligning the robot's internal odometry-based yaw to the ground truth provided by Gazebo. The script corrects whatever accumulated heading drift the odometry has already developed by the time the calibration commences, snapping it to the nearest cardinal direction, rather than attempting to cancel it out entirely, and allowing the in-place rotation to make use of a corrected heading.

The node tracks three related but distinct yaw angles: the raw value, \texttt{raw\_yaw}, obtained by converting the \texttt{/odom} orientation quaternion into degrees, the yaw offset, \texttt{yaw\_offset}, and the current, offset-corrected yaw, \texttt{current\_yaw}, used to control the in-place rotation. This separation allows the calibration to take the accumulated drift into account while still ensuring that once the correction gets applied, the operator-facing yaw value snaps to the nearest cardinal direction, rather than retaining any fractional degree value. The separation of the raw and current yaw also facilitates the calibration process itself; as the node runs, the \texttt{raw\_yaw} value continues to update to reflect the robot's actual heading, while the offset remains constant, thus adjusting the in-place rotation to match the corrected value.

The calibration process then involves querying the ground-truth orientation of the robot, obtaining it directly from Gazebo by calling \texttt{gz topic -e -t /world/world\_demo/pose/info --json-output -n 1} as a subprocess and parsing the JSON string to extract the relevant quaternion. The tool then uses this value to calculate the offset required to align the odometry to the ground-truth heading, subtracting the offset from the \texttt{raw\_yaw} to obtain the \texttt{current\_yaw} value. This approach to calibration deliberately takes advantage of a feature only available in the simulated setup, namely, the ability to access the ground truth, rather than attempting to estimate it based on available sensors.

The yaw offset in turn gets calculated via the rounded gazebo\_deg value, by using \texttt{round(gazebo\_deg / 90.0) * 90.0}, rather than applying it directly as the correction. Calibration then snaps the \texttt{current\_yaw} to the nearest cardinal axis, rather than attempting to retain the exact reading from the ground-truth orientation. The reasoning behind this lies in the practicality of visual inspection; using the SLAM toolchain to build a map of the warehouse environment entails building a scan-matching trajectory, which in turn gets converted into a 2D occupancy grid~\cite{macenski2021slam}. As this grid gets visualized in Rviz, having the long axis of the warehouse map align with the row and column axes of the display makes it significantly simpler to interpret, inspect, and reason about the generated map, relative to an identical map rotated by an arbitrary value~\cite{ribeiro2025indoor}. In practice, the calibration only needs to provide a sufficiently accurate coarse alignment to facilitate visual inspection; the resulting error from using a cardinal-axis snap of up to $45^\circ$ to the true heading is insignificant at the scale of a warehouse environment.

A similar fallback is available for cases where Gazebo calibration becomes unavailable, either due to the calibration subprocess exiting early or the calibration topic being unreachable. In such cases, when a call to the calibration topic fails, the script snaps the \texttt{raw\_yaw} to the nearest cardinal axis and warns the user that the ground-truth calibration attempt has failed. This lets the user retain the convenience of cardinal snapping, while being warned that the alignment is no longer accurate. The approach represents a variation of the theme of retaining partial functionality under adverse conditions, and mirrors the behavior of the other calibration and diagnostic tools described in this section.~\cite{kagizman2026lowcost}.

\section{Operator Interface and Human Supervisory Control}

\subsection{Thread-Safe PyQt5-ROS 2 Bridging}

\texttt{monitoring\_dashboard.py} ties together two event loops that were never designed to run in the same thread: Qt's, which must run on the main thread to allow for safe widget manipulation, and \texttt{rclpy}'s, which must constantly spin to process subscriptions. The most straightforward way to bind the two, which would involve directly calling a Qt widget update method from within a subscription callback, is unsafe to do regardless of how light the called update is, because Qt widgets are not reentrant and a callback triggered by \texttt{rclpy}'s executor is not guaranteed to run on the main thread anyway.

Instead, the dashboard application uses the signal-slot mechanism to communicate between the two threads, which is safe to use directly: a signal emitted from any thread will be delivered to the slot connected to it on the receiving thread's event loop, which for a Qt main window is the main thread. \texttt{ROS2SpinThread}, a simple \texttt{QThread}-based class, is used to isolate the ROS~2 executor loop on a separate thread; its run method simply calls \texttt{rclpy.spin(self.node)}, which puts the node into the continuous spinning mode. The \texttt{DashboardROS2Node} node's subscription callbacks are not written to directly manipulate any widgets; instead, each callback parses the received message and emits a signal declared on \texttt{DashboardBridge}, a \texttt{QObject}-based class shared with the GUI window. The latter connects each of the five signals, declared to carry status, camera frames, anomaly alerts, shelf arrival notifications, and mission-armed status, to its respective update slot in \texttt{\_connect\_signals}; each widget update is then safely performed on the main thread, regardless of which thread the signal was emitted from.

The same signal-slot mechanism imposes a discipline on the \texttt{DashboardROS2Node}'s callback functions: while they are free to parse and validate the contents of the received messages, they are not permitted to directly update any widgets, as that would involve accessing objects that reside in the main thread's memory space from a different thread. The camera-frame callback demonstrates a case where this discipline requires extra care even when the callback does not update any widgets directly. \texttt{\_camera\_callback} reconstructs a \texttt{QImage} from the incoming \texttt{sensor\_msgs/Image} message's data buffer, but makes a deep copy of it before emitting the image via the signal, rather than using the buffer directly. A \texttt{QImage} constructed from a buffer pointing to external memory, as one would get from the Image message's data directly, is a view onto that buffer; if the emitting callback returns to the main thread and the Image message it was received with is deallocated, or if a subsequent frame overwrites that buffer before the GUI thread reads it, the \texttt{QImage} in the GUI thread will end up pointing to invalid memory. This can manifest as intermittent segmentation faults in the dashboard application, depending on whether the memory the buffer pointed to was reused by other ROS~2 messages or simply unmapped, and is notoriously difficult to debug due to the non-deterministic nature of the fault. Moving the buffer contents to a separate \texttt{QImage} with a deep copy adds a small overhead per frame but avoids the hazard entirely, as the buffer now resides in a memory block allocated specifically for the \texttt{QImage} and will be deallocated when that object is destroyed~\cite{blass2022real_thesis}.

A second subtlety in the dashboard application's thread management comes at the process level: the \texttt{main} function sets \texttt{signal\_handler\_options=SignalHandlerOptions.NO} when initializing \texttt{rclpy} and sets Python's default SIGINT handler to be invoked, rather than the one installed by Qt. If left at their defaults, both \texttt{rclpy} and Qt's event loop would attempt to handle SIGINT by initiating their respective shutdown procedures, but doing so from a signal handler is unsafe in the first place and the two handlers could interfere with one another, potentially leaving the process in an inconsistent state. Instead, by restoring the default handler and relying on \texttt{closeEvent} to handle SIGINT, the dashboard application moves the shutdown logic to the main thread: \texttt{closeEvent} itself is called when the GUI window is closed, but a SIGINT sent to the process will also reach the default handler and trigger the same sequence. The shutdown procedure now cleanly terminates the ROS~2 node with \texttt{rclpy.try\_shutdown()}, which is called both from the \texttt{closeEvent} handler and, explicitly, at the end of the \texttt{main} function after \texttt{app.exec\_()} returns.

The signal-based bridge established between the ROS~2 node and the Qt main window carries specific information about the content of the detected anomalies. When the consolidated \texttt{waypoint\_sender} node's scanning pipeline classifies a section of the environment as damaged, it projects the attribution heatmap from the backprojection in Equation~\ref{eq:xai_back_projection} onto the raw sensor frame, blending the two images together into a single \texttt{sensor\_msgs/Image} message published on the robot's \texttt{camera\_view} topic; a similar procedure is used for intact sections, but with no heatmap overlay. The dashboard's \texttt{\_camera\_callback} discriminates between the two cases based on the image message's encoding: a three-channel \texttt{bgr8} buffer is converted to \texttt{QImage}'s \texttt{QImage.Format\_RGB888} format and then passed to \texttt{rgbSwapped()} to account for the difference in channel order between OpenCV and Qt, while a one-channel \texttt{mono8} buffer is converted to \texttt{Format\_Grayscale8} directly. The same considerations apply to the \texttt{QImage}s constructed from the buffer in both cases, with the deep copy discussed earlier being performed before the image is emitted as a signal; the channel-order correction and the memory-safety correction are thus both performed at the same point in the code same point in the code~\cite{macenski2022robot,albatati2026ros2}.

\subsection{Feedback-Driven Window Tiling under WSLg}

The three primary visual displays used by the operator throughout a mission run - RViz2, the Gazebo GUI client, and the PyQt5 monitoring dashboard of Section~7.6.1 - each run as independent X11 windows under WSLg, and are not positioned or sized in any particular relation to each other by default. \texttt{window\_tiler.py} exists to lay out these three windows in a particular set of non-overlapping positions and sizes, so that the operator may view them all simultaneously without manual intervention.

The simplest possible approach to this task would be to issue a single \texttt{xdotool windowmove} and \texttt{window\-size} command pair to each target, and declare the operation complete when they each return successfully. \texttt{tile\_window\_stabilized}, however, performs additional steps after each command in order to work around limitations in WSLg's XWayland composition stack. It queries \texttt{xdotool getwindowgeometry} on the target after each movement/resizing operation, and compares the resulting position and size to the intended one with a tolerance of 15 pixels per dimension, exiting early if the observed geometry matches the target. This is necessary because the underlying Windows desktop window manager and its interaction with WSLg's compositor do not always respect a given window size or position - the WSLg documentation does not indicate that applications should expect any particular behavior beyond what is mandated by the X11 protocol. Consequently, blindly issuing resize/move commands and assuming them to have taken effect would cause windows to appear in unplanned positions under certain circumstances, with no feedback as to what had occurred.

Issuing the same resize/move sequence repeatedly at this point would introduce another potential issue. If a window has an invisible pixel-wide border of configurable size that the WSLg window manager places around every top-level window which the \texttt{xdotool} geometry queries do not account for, then such a window placed at (0,0) with a requested size of (800, 600) would have an observed position of (0,0) and size of (800 - $dx$, 600 - $dy$), where $dx$ and $dy$ are the thicknesses of the aforementioned borders. Further retries of the same resize operation would observe the same $dx$ and $dy$ and conclude erroneously that the window manager had failed to honor the resize completely, when in reality, the window had been placed successfully against an invisible border that is not honored by the geometry query. A naive retry loop here would enter an infinite loop trying to adjust the window position, or introduce a hard-coded retry limit that would leave the window placed at the position that had failed, rather than its intended position.

The function \texttt{tile\_window\_stabilized} avoids this situation by comparing the observed $dx$ and $dy$ against the values recorded on the immediately preceding attempt, on every attempt from the second onward, and exiting early the first time the two agree. Rather than considering repeated failures to match the requested geometry to be a sign of an obstructing border, it observes that the observed $dx$ and $dy$ are stable and considers the window to have been placed successfully to a position which the WSLg window manager has explicitly stated to be invalid for top-level windows.

This allows it to avoid an unacceptable false positive rate while still honoring the observed behavior of the WSLg window manager. It is not feasible to distinguish between a transient system-level issue preventing a resize/move from taking effect, and a stable but erroneous placement at a position that the system explicitly forbids for top-level windows, but it is still beneficial to avoid treating the latter as a success condition when the behavior of the window manager is known to be stable across multiple attempts. By distinguishing between a placement that matches the target dimensions, one that is stable but offset, and one that mismatches the target dimensions with no stable offset, the function makes it possible for a retry loop to distinguish between these three cases and terminate when it has observed behavior that matches expected system behavior at the cost of complexity in its retry logic and an unacceptable false-positive rate. This places a five-cycle limit on the number of resize/move operations that can be attempted before the function proceeds to the next window, and limits the amount of time between each attempt to $0.15\,\text{s}$, in order to allow WSLg's underlying RDP-based window compositing to stabilize before querying \texttt{xdotool} again for the updated geometry~\cite{microsoft2024wslg}.

The main loop surrounding this retry logic operates within a higher-level bounded loop, with a $120\,\text{s}$ hard-coded timeout. It detects the presence of each of the three target windows using \texttt{xdotool search}, discarding any windows smaller than $300\times300$ pixels as false positives (such as splash screens that appear briefly during RViz2 and Gazebo startup), and waits for each of the three target windows to appear in the system before proceeding to tile any of them to their intended positions and sizes. It is not possible for a single window to cause the main loop to fail entirely, since its retry loop operates separately from the discovery logic - and while RViz2, being the last launched of the three windows as per the launch configuration of Section~7.2.3, may cause the tiling of the other two windows to be deferred until after it has started, each of the three windows is tiled independently as soon as it becomes available. The timeout exists only to avoid spending an unbounded amount of time waiting for one of the windows to appear when it never does.

\section{Simulation Substrate and Mapping Workflow}

\subsection{Gazebo Harmonic Rendering Configuration for Headless-Adjacent Operation}

Gazebo Harmonic's default GUI configuration loads a comparatively large set of plugins, an outliner tree, a plugin browser drawer, an entity inspector, plus the 3D viewport itself, an appropriate set for interactive scene authoring but a bewildering set of distractions for a simulation whose 3D view is a passive monitor for the operator, who is actually working in the PyQt5 dashboard and the RViz2 windows tiled alongside it, as described in Section~7.6.2. \texttt{gz\_gui.txt} specifies a much minimized GUI configuration, which retains only the three plugins necessary to actually see anything in a 3D view, \texttt{MinimalScene} for rendering, \texttt{GzSceneManager} as the backing scene-graph manager, and \texttt{InteractiveViewControl} for mouse panning and zooming~\cite{soriano2026benchmarking}, with the drawer menu and the plugin browser explicitly disabled via \texttt{visible="false"}.

The two backing plugins, \texttt{GzSceneManager} and \texttt{InteractiveViewControl}, are in turn loaded in a floating window constrained to $5\times5$ pixels with the title bar suppressed, rather than being omitted from the configuration entirely; this suggests a particular dependency on those two plugins, which are necessary even if not visible, by the rendering plugin \texttt{MinimalScene} -- camera and lighting setup are controlled by \texttt{GzSceneManager} rather than done inline by \texttt{MinimalScene}, which in turn has been noted to cause renderer initialization failures under the Ogre2 backend if either of those two plugins is not loaded at all, with the note in the configuration itself documenting this as the rationale for not simply omitting them. The camera itself is fixed to a specific top-down pose, \texttt{camera\_pose} explicitly set to a position just above the center of the warehouse floor and looking straight down rather than at whatever default angle a fresh Gazebo session would have used. This is actually preferable for the limited role of the 3D view in this particular setup, which is to provide general situational awareness of the three robots operating within the warehouse; a fixed top-down view is substantially more immediately informative than an arbitrary perspective view, especially considering that the same warehouse layout is already being shown in a top-down projection within RViz2's own costmap and path displays. Fixing the camera pose within the configuration itself rather than leaving it to the operator to re-center every time the 3D view is opened reduces yet another potential source of human error, particularly when compared to the deterministic launch orchestration described in Section~7.2, which this configuration's minimized set of plugins is largely a prerequisite for.

This configuration thus represents a compromise between a fully headless simulation, which would forgo any 3D rendering entirely and save on GPU resources, and the fully-featured default GUI configuration. The former would remove any possibility for visual confirmation that the physics engine is operating as intended, with the risk of silent simulation errors becoming particularly acute if Gazebo's default physics engine is used and the operator is relying on the dashboard and RViz2 for real-time telemetry rather than any form of video or photonic capture of the actual screen content; it becomes substantially easier to spot if a robot has fallen through the floor by looking at a live 3D rendering than it is to guess at the underlying cause of a given localization error based on ROS~2 topic data. Leaving a minimal GUI client running with a static camera pose avoids that risk while imposing minimal additional overhead on the system, aside from the CPU and RAM cost of maintaining a separate Gazebo client process, which is already necessary for the actual autonomous operation of the robots~\cite{koenig2004design}.

\subsection{Cartographer-Based Offline Map Generation}

The occupancy grid map used by all three robots' \texttt{map\_server} instances under autonomous operation, as well as extensively referenced in Chapter~4 when discussing the physical workspace $\mathcal{W}$, is not generated by hand but rather results from a separate single-robot mission run in the same warehouse, launched via \texttt{slam.launch.py} and configured in \texttt{warehouse\_cartographer.lua}. This procedure is deliberately decoupled from the three-robot autonomous operation stack described in Chapters~5 and~6; it instantiates a single, unnamespaced TurtleBot3 instance in the warehouse world and drives it around via the interactive yaw calibration tool described in Section~7.5.2 until a satisfactory map is generated by Cartographer's occupancy-grid node, before terminating the process early to avoid launching any of the Nav2 stacks, claim management protocols, or task allocation mechanisms used in later chapters~\cite{macenski2020marathon}.

The configuration for Cartographer's SLAM stack itself sets \texttt{provide\_odom\_frame} to false and \texttt{\seqsplit{published\_frame}} to \texttt{odom} rather than having it publish its own internal odometry-rooted transform, in a similar vein to the dual TF publisher resolution used for the production EKF in Section~7.3.2. Cartographer is itself a potential second writer to the \texttt{map} $\rightarrow$ \texttt{odom} edge, since its trajectory builder maintains its own local estimate of the robot's position independent of whatever odometry is fed to it, and publishing that estimate as a competing transform would reproduce the same conflict the EKF configuration in Section~7.3.2 was built to eliminate. By configuring \texttt{provide\_odom\_frame} to false and \texttt{use\_odometry} to true, Cartographer is instructed to consume the externally-supplied odometry input rather than publish a competing \texttt{/tf} tree of its own, keeping the single-writer discipline established elsewhere in this chapter consistent across both the offline mapping pass and the autonomous mission configuration, even though the mapping pass, as already noted, never runs alongside any Nav2 stack of its own.

The $20.0\,\text{m}$ laser range configured for the front and rear depth sensors in use for this deployment, and the reasoning behind it as described in Chapter~4, Section~4.2 for allowing the particle filter to observe the central structural pillars of the warehouse, apply just as well for this standalone mapping mission as they do for the autonomous robots. It is notable that the same sensor configuration is used both for Cartographer's scan matcher and for AMCL for the actual autonomous navigation, although the same could be said for the use of the same particle filter configuration across all three \texttt{map\_server} instances in use for the three robots.

The loop closure parameters for Cartographer's \texttt{optimize\_every\_n\_nodes} of 120 and the constraint-builder minimum score threshold of 0.75 dictate how aggressively the SLAM algorithm will attempt to correct its position estimates once a loop closure has been detected, which is relevant to the mapping procedure itself in this particular warehouse due to the narrow aisle geometry described in Chapter~4, Section~4.1.2. The mapping trajectory, manually driven between pairs of shelves, accumulates localization drift over the length of an aisle due to the lack of visual and geometric landmarks along the length of the warehouse, with the narrow aisle width limiting the correction maneuvers available to the robot~\cite{ribeiro2025indoor} -- it is at this point that the pose graph's loop closure algorithm is substantially more effective than attempting to cancel the drift with an opposing movement along the same axis, the drift along the global X axis being canceled by a movement along the global negative Y axis once a loop has been completed~\cite{hess2016real}.

Since this mapping mission serves only to generate a static reference map, rather than needing to dynamically adapt to the changing warehouse state the way the autonomous mission stack of Chapters~5 and~6 must, the process only needs to be run once against any given warehouse configuration rather than every time the autonomous robots are launched. The pre-flight stale-process cleanup procedure described in Section~7.2.4 is a direct result of this particular requirement, as an interrupted mapping mission could leave a hung Cartographer process from the previous attempt interfering with the fresh mapping process, thereby corrupting the newly-generated occupancy grid that the autonomous robots will rely on for the duration of their own mission.

\section{Integration Summary and Transition to Experimental Validation}

The engineering contributions detailed across this chapter serve to realize the behavioral models of Chapters~5 and~6 on an actual physical WSL2 host as a process which may be subjected to empirical observation and recalibration. Four themes are revealed to which solutions across seemingly disconnected challenges in this domain are found to conform. First, in the consolidated node graph of Section~7.1, the sole-TF-writer mechanism of Section~7.3.2, and the absolute-path topic exception of Section~7.3.3, a theme emerges of responsibility for a shared resource or authority being re-parented to a single entity to resolve potential ambiguity in favor of an assertion by one party over the arbitration by time that the shared resource or authority would otherwise undergo. Second, in the gate-based launch sequencing of Section~7.2, readiness for a downstream process to begin execution is made conditional on a polled observation of an externally visible state, rather than an assumption of sufficient elapse since a previous event as is practiced in the diagnostic tools of Section~7.4. Thirdly, the fallback calibration sequence of Section~7.5.2 and the stable-drift acceptance criteria of Section~7.6.2 each represent a tolerance for partial failure in the form of a lesser alternative beginning to take effect when a preferred alternative is unavailable, a pattern which also finds incorporation at the design level in Chapter~6. Finally, a separation of concerns is observed between the low-level raw data and higher-level derived data in both the CSV versus summary visualization of Section~7.4.4 and the superseded-but-retained relay mechanism of Section~7.5.1.

Lastly, it is noted that none of the contributions of this chapter change the fundamental mathematical properties of the behavioral models introduced in Chapters~5 and~6, but rather serve to ensure that these theories can be subjected to empirical evaluation by actually being built and run on real hardware. The experimental campaigns of Chapter~8 will depend on this engineering infrastructure to actually take place, as all quantitative results and observations are inherently reliant on the correctness of the underlying telemetry and instrumentation.